% Section 15.8, from lectures/L15/appendix/b_exercises.md.

\section{Uppgifter}
\label{c15:sec:uppgifter}

\begin{aside}
\textbf{Så arbetar du med uppgifterna.} Del~1 och del~2 räknas under lektionen: först på egen
hand, några minuter per uppgift, sedan gemensamt. Del~3 är extrauppgifter att träna vidare på
efter lektionen. De flesta går att räkna i huvudet eller med papper och penna.

\facitnotis{L15}
\end{aside}

% ------------------------------------------------------------------------------------------
\uppgiftsdel{Del 1}{Repetitionsuppgifter}

\uppgift{1.1}{Återfå tidsfunktion från derivata}
En sensorspänning $u(t)$ i en mätkrets är inte känd direkt, men dess derivata är uppmätt till
\[
  u'(t) = (6t - 4)\,\text{V/s}.
\]
Det är också känt att startspänningen är $u(0) = 2\,\text{V}$.
\begin{parts}
  \item Bestäm ett uttryck för $u(t)$.
  \item Bestäm spänningen efter $3$ sekunder.
\end{parts}

% ------------------------------------------------------------------------------------------
\uppgiftsdel{Del 2}{Nytt stoff}

\uppgift{2.1}{Rektangulär $\rightarrow$ polär form}
Omvandla spänningen $U = 5 - j2\,\text{V}$ till polär form.

\uppgift{2.2}{Polär $\rightarrow$ rektangulär form}
Omvandla strömmen $I = 5\angle\frac{\pi}{6}\,\text{mA}$ till rektangulär form.

\uppgift{2.3}{Beräkning av impedans}
En krets matas med spänningen från \uppref{2.1}, och strömmen från \uppref{2.2} flyter genom
kretsen. Beräkna kretsens impedans $Z$ med ”Ohms lag”:
\[
  Z = \frac{U}{I},
\]
där
\begin{itemize}
  \item $Z$ är kretsens impedans i $\Omega$,
  \item $U$ är matningsspänningen i V,
  \item $I$ är strömmen genom kretsen i mA.
\end{itemize}
Svara både på polär och på rektangulär form.

% ------------------------------------------------------------------------------------------
\uppgiftsdel{Del 3}{Extrauppgifter}
Uppgifterna nedan är extra träning och görs med fördel på egen hand efter lektionen.

\uppgift{3.1}{Den imaginära enheten}
Beräkna.
\begin{delar}(5)
  \task $j^2$
  \task $j^3$
  \task $j^4$
  \task $j^5$
  \task $\dfrac{1}{j}$
  \task* Lös ekvationen $x^2 = -25$.
\end{delar}

\uppgift{3.2}{Real- och imaginärdel}
De komplexa talen $z_1 = 4 - j3$, $z_2 = -2 + j5$, $z_3 = j6$ och $z_4 = -7$ är givna.
\begin{parts}
  \item Ange $\operatorname{Re}(z)$ och $\operatorname{Im}(z)$ för samtliga tal.
  \item Vilket tal är rent imaginärt?
  \item Vilket tal är rent reellt?
\end{parts}

\uppgift{3.3}{Addition och subtraktion}
Beräkna med $z_1 = 3 + j2$ och $z_2 = -1 + j5$:
\begin{delar}(4)
  \task $z_1 + z_2$
  \task $z_1 - z_2$
  \task $2z_1 + 3z_2$
  \task $z_2 - 2z_1$
\end{delar}

\uppgift{3.4}{Multiplikation}
Beräkna och skriv svaret på rektangulär form.
\begin{delar}(4)
  \task $(2 + j3)(1 + j4)$
  \task $(3 - j2)(3 + j2)$
  \task $j(5 - j2)$
  \task $(1 + j)^2$
\end{delar}

\uppgift{3.5}{Konjugat}
\begin{parts}
  \item Ange konjugatet till $5 + j2$.
  \item Ange konjugatet till $-3 - j7$.
  \item Beräkna $z\bar{z}$ för $z = 4 + j3$.
  \item Visa allmänt att $z\bar{z} = \abs{z}^2$.
\end{parts}

\uppgift{3.6}{Division}
Beräkna och skriv svaret på rektangulär form.
\begin{delar}(3)
  \task $\dfrac{4 + j2}{1 - j}$
  \task $\dfrac{6 + j8}{j}$
  \task $\dfrac{10}{3 + j4}$
  \task* Varför multiplicerar man med nämnarens konjugat vid division?
\end{delar}

\uppgift{3.7}{Absolutbelopp}
Beräkna.
\begin{delar}(5)
  \task $\abs{3 + j4}$
  \task $\abs{-5 + j12}$
  \task $\abs{j7}$
  \task $\abs{-8}$
  \task $\abs{1 - j}$
\end{delar}

\uppgift{3.8}{Fasvinkel med kvadrantkorrigering}
Bestäm fasvinkeln. Ange vilken kvadrant talet ligger i.
\begin{delar}(5)
  \task $4 + j4$
  \task $-4 + j4$
  \task $-4 - j4$
  \task $4 - j4$
  \task $3 + j4$
\end{delar}

\uppgift{3.9}{Rektangulär till polär form}
Skriv på polär form.
\begin{delar}(4)
  \task $z = 6 + j8$
  \task $z = -3 + j4$
  \task $z = -1 - j1$
  \task $z = -j5$
\end{delar}

\uppgift{3.10}{Polär till rektangulär form}
Skriv på rektangulär form.
\begin{delar}(3)
  \task $z = 10\angle 30^{\circ}$
  \task $z = 4\angle 90^{\circ}$
  \task $z = 6\angle 180^{\circ}$
  \task $z = 2\angle{-60^{\circ}}$
  \task $z = 8\angle\dfrac{\pi}{4}$
\end{delar}

\uppgift{3.11}{Spänning som komplext tal}
En spänning ges av $U = 12 + j5\,\text{V}$.
\begin{parts}
  \item Beräkna $\abs{U}$.
  \item Beräkna fasvinkeln $\delta$.
  \item Skriv $U$ på polär form.
  \item Vad representerar fasvinkeln fysikaliskt?
\end{parts}

\uppgift{3.12}{Impedans på rektangulär form}
En seriekrets har $Z_R = 30\,\Omega$ och $Z_L = j40\,\Omega$.
\begin{parts}
  \item Bestäm den totala impedansen $Z$ på rektangulär form.
  \item Beräkna $\abs{Z}$.
  \item Beräkna impedansens fasvinkel.
  \item Är kretsen induktiv eller kapacitiv?
\end{parts}

\uppgift{3.13}{Ohms lag med komplexa tal}
En krets med impedansen $Z = 3 + j4\,\Omega$ matas med spänningen $U = 10\,\text{V}$ (rent reell).
\begin{parts}
  \item Beräkna $\abs{U}$ och $\abs{Z}$.
  \item Beräkna strömmen $I = \dfrac{U}{Z}$ på rektangulär form.
  \item Beräkna $\abs{I}$ och kontrollera mot $\dfrac{\abs{U}}{\abs{Z}}$.
  \item Beräkna strömmens fasvinkel och tolka resultatet.
\end{parts}

\uppgift{3.14}{Andragradsekvationer med komplexa rötter}
Lös ekvationerna.
\begin{delar}(3)
  \task $x^2 + 9 = 0$
  \task $x^2 + 4x + 13 = 0$
  \task $x^2 - 2x + 5 = 0$
  \task* Vad gäller alltid för de två rötterna när diskriminanten är negativ?
\end{delar}

\uppgift{3.15}{Serie- och parallellimpedans}
Två impedanser $Z_1 = 4 + j3\,\Omega$ och $Z_2 = 4 - j3\,\Omega$ är givna.
\begin{parts}
  \item Beräkna seriekopplingen $Z_1 + Z_2$ och kommentera resultatet.
  \item Beräkna $\abs{Z_1}$ och $\abs{Z_2}$.
  \item Beräkna produkten $Z_1Z_2$.
  \item Beräkna parallellkopplingen $\dfrac{Z_1Z_2}{Z_1 + Z_2}$.
\end{parts}

\uppgift{3.16}{Sant eller falskt}
Avgör om påståendet är sant eller falskt och motivera kortfattat.
\begin{parts}
  \item $j^2 = 1$
  \item Absolutbeloppet $\abs{z}$ kan vara negativt.
  \item Ett rent imaginärt tal har fasvinkeln $90^{\circ}$ eller $-90^{\circ}$.
  \item Produkten $(a + jb)(a - jb)$ är alltid reell.
  \item Räknarens $\arctan$ ger alltid rätt fasvinkel direkt.
  \item Addition av komplexa tal är enklast på polär form.
\end{parts}
